\section{Reviewer 3}

\subsection{Major Comments}

\begin{enumerate}
	\item It is not entirely clear at any point in this manuscript what problem the manuscript is trying to address. The author(s) are correct that civil war scholars have not properly considered the interdependencies among actors in multiparty conflicts and that dyadic research designs are limited. However, what are the consequences of these limitations? What don't we understand that, at the end of this article, we will understand better?
	\begin{itemize}
		\item \textcolor{blue}{ \emph{
		By re-writing the introduction, theory, and concluding sections, we hope that the main point of the paper is much clearer. We focus on the question of ``Who fights who and when" in the context of multi-group armed conflicts. We argue, and find support, for the position that the interactions, or behavior, of armed actors are just as important for explaining future occurrences of violence as armed group attributes. Throughout the paper, but especially in the introduction, we suggest that by missing the interactions present in violent conflict scholars will mischaracterize patterns of violence and ignore extra-dyadic factors that make violence between groups more likely.
		}}
	\end{itemize}	
	\item The author(s) state on page 3 "This study investigates how relationships between armed actors evolve over time, and how the structure of relationships in one time period influences the occurrence of violence between armed actors in a future period. To do this, we consider how interdependence across actors, changes in actor composition over time, and actor-level attributes predict violence in the network." There are several potential questions here, such as whether different armed groups fight (or cooperate), how the presence of additional actors affects the level of violence between rebel groups and the state, the overall level of violence in the conflict, etc. The lack of specificity comes through in the theory, where there is a discussion of key concepts in network theory like actor centrality and reciprocity, but it is unclear how exactly these concepts relate to the question in the manuscript or to existing work. It is also present in the analysis, where the DV is the occurrence of battles between armed actors, but it is not clear why that is appropriate given the theory.	
	\begin{itemize}
		\item \textcolor{blue}{ \emph{
		We have now re-written the introduction and the theory sections. First, we hope that our introduction is much clearer and orients the paper around the question ``who fights who and when"? We illuminate this puzzle through providing anecdotes about different multi-actor conflicts in our new introduction. Our theory section is now rooted in contributions from prior research on the attributes and motivations of armed actors. Such a focus allows us to then explain how a network approach flows from the findings of preexisting research. Each network concept used in the theory is now situated within the related literature on conflict processes, which allows us to defend our position that the interactions between actors, not just actor attributes, more accurately explains the occurrence of violence between armed groups. 
		}}
	\end{itemize}	
	\item All of this is fixable, but it requires re-framing the introduction, better situating the argument in the literature, and re-writing the theory. Why not start with the question� why do some armed groups fight each other and others do not? There is clear literature on this question, much (although not all) of which is cited here. Describe that literature and explain why it is insufficient to understand this phenomenon. Then, the manuscript needs a clear, focused theoretical argument drawing on network theory to explain this. The theory should start from first principles� what are armed actors trying to achieve through violence and how do the components here (centrality, reciprocity, the introduction of new ``aggressive'' actors, etc.) fit into this story.	
	\begin{itemize}
		\item \textcolor{blue}{ \emph{
		We agree with this feedback and, as noted above, have completely rewritten the front end of the paper in an effort to clarify our focus on this question of why some groups fight each other.  We have also retitled the paper to reflect this focus. 
		}}
	\end{itemize}		
\end{enumerate}

\subsection{Minor Comments}

\begin{enumerate}
	\item To identify groups, the author(s) use ACLED. I understand why, but the downside of doing this is that they lose the ability to include group level characteristics present in, for example, the Non-State Actor data, the EPR project, etc. The authors need to make a clearer case for the contribution here given that they are giving up substantial access to information about these actors.
	\begin{itemize}
		\item \textcolor{blue}{ \emph{
		We make the case for our use of the ACLED data on page 14, footnote 11. Our general motivation for using the ACLED data is that more actors are included in the data because there is no threshold requirement, such as a specific number of battle-related deaths, like there is in the UCDP data. This is useful for a study interested in explaining fighting between actors since not all battles will result in a high number of deaths, but all battles might influence the trajectory of the conflict.
		}}
	\end{itemize}	
	\item The manuscript needs to be much clearer on what the system is. Here, it includes everyone fighting in Nigeria. But is that always appropriate? Is the country the box that bounds networks? Why? What about separate, peripheral, conflicts in places like India, should they be modelled together? What about linked conflicts in places like Central Africa? Should they be modelled separately? These considerations are important.
	\begin{itemize}
		\item \textcolor{blue}{ \emph{
		We discuss how the modeling framework would differ given transnational or subnational conflicts on page XXX.
		}}
	\end{itemize}	
	\item The manuscript should, at the least, cite Bapat and Bond (2011), which is one of the earliest articles to examine cooperation/conflict between rebel groups in civil war.
	\begin{itemize}
		\item \textcolor{blue}{ \emph{
		We thank you for pointing out this oversight. We now discuss Bapat and Bond's paper on page 7.
		}}
	\end{itemize}				
\end{enumerate}